\section{Experiments}

\begin{figure*}[!t]
    \centering
    \begin{subfigure}{0.33\linewidth}
        \centering
        \includegraphics[width=\linewidth]{fig/RCE.pdf}
        \caption{Remove Code Execution}
        \label{fig:sub1}
    \end{subfigure}%
    \begin{subfigure}{0.33\linewidth}
        \centering
        \includegraphics[width=\linewidth]{fig/RCE.pdf}
        \caption{Sensitive Information Disclosure}
        \label{fig:sub2}
    \end{subfigure}%
     \begin{subfigure}{0.33\linewidth}
        \centering
        \includegraphics[width=\linewidth]{fig/RCE.pdf}
        \caption{Insecure File Operation}
        \label{fig:sub3}
    \end{subfigure}%
    \hfill

     \begin{subfigure}{0.33\linewidth}
        \centering
        \includegraphics[width=\linewidth]{fig/RCE.pdf}
        \caption{SQL Injection}
        \label{fig:sub4}
    \end{subfigure}%
     \begin{subfigure}{0.33\linewidth}
        \centering
        \includegraphics[width=\linewidth]{fig/RCE.pdf}
        \caption{Denial of Service}
        \label{fig:sub5}
    \end{subfigure}
    \vspace{-10pt}
    \caption{Experimental Results w.r.t MCP Tool Vulnerability Types}
    \label{fig:main}
\end{figure*}


\begin{table}[!t]
\centering
\small\caption{Experiment Results}
\vspace{-10pt}
    \begin{tabular}{m{2cm}m{1.5cm}m{1.5cm}m{1cm}m{1cm}}
\toprule
\textbf{Model} & \textbf{Precision} & \textbf{Recall} & \textbf{F1-Score} & \textbf{MPA} \\
\midrule
GPT-5-mini & XX.X\% & XX.X\% & XX.X\% & XX.X\% \\
Claude Haiku-4-5 & XX.X\% & XX.X\% & XX.X\% & XX.X\% \\
DeepSeek & XX.X\% & XX.X\% & XX.X\% & XX.X\%\\
AI-Infra-Guard & XX.X\% & XX.X\% & XX.X\% & XX.X\%\\
Bandit & XX.X\% & XX.X\% & XX.X\% & XX.X\%\\

\bottomrule
\end{tabular}

\label{tab:eva}
\end{table}




\subsection{Experiment Setup}

\textbf{Evaluated Baselines.} We evaluate a diverse set of baselines, ranging from state-of-the-art LLMs to industry-standard security analysis platforms, including \textbf{GPT-5-mini}~\cite{gpt5mini}, \textbf{Claude Haiku 4.5}~\cite{claude}, \textbf{DeepSeek-V3.2}~\cite{DeepSeekAI2025DeepSeekV32PT} , \textbf{AI-Infra-Guard}~\cite{aiinfraguard} (Released by Tencent Zhuque Lab in 2025) and \textbf{Bandit}~\cite{badit} (Originally developed by the OpenStack Security Team).

\textbf{Configuration.} We set the generation temperature to 0.0, disable streaming, and use a 120-second timeout per API request. The evaluation employs dual prompts: a system prompt defining the five MCP vulnerability types, and a user prompt containing the raw Python code of the target server. Security tools (Bandit and AI-Infra-Guard) are run with default settings. To prevent models from exploiting sample order, we shuffle the JSONL dataset with \texttt{random.shuffle()} and a fixed seed, ensuring each sample is evaluated independently. To ensure experimental fairness, we disabled the dynamic penetration testing components of AI-Infra-Guard, and retained only its static analysis functionality in our experimental setting.



\textbf{Prompt Design.} We use a zero-shot, role-playing prompting strategy for LLM evaluation, separating the prompt into a system and a user prompt to ensure clarity and output stability. The system prompt assigns the role of an “expert security researcher” for Python MCP auditing and embeds the full vulnerability taxonomy, including code-level indicators across five categories. Responses are strictly constrained to a JSON schema with \texttt{is\_vulnerable}, \texttt{vulnerability\_type}, and \texttt{reasoning} fields. The user prompt provides only the raw Python source code.


\subsection{Results}



\todo{Table \ref{tab:eva2} presents xx.}
